\section{传统供需基准模型}

\subsection{模型建立}

传统供需基准模型用于刻画“没有显著缓冲机制时”的理论价格压力。它不是最终预测模型，而是反事实参照：如果市场只按低短期弹性吸收供应中断，那么静态供需框架会给出多高的价格压力。

本文主口径采用线性需求曲线假设。设局部反需求函数为 \(P=a-bQ\)，并在战前均衡点 \(P_0,D_0\) 处使用题面给出的短期点弹性 \(\varepsilon\) 标定斜率。当供应中断量为 \(I\) 时，线性需求反事实价格为

\begin{equation}
  P_{\mathrm{base}}
  =
  P_0
  \left(
    1+\frac{I/D_0}{|\varepsilon|}
  \right).
\end{equation}

需要说明的是，上式不是常弹性需求函数的一阶 Taylor 近似，而是线性需求反事实基准的主口径。选择线性需求，是因为它隐含消费者边际支付意愿随消费量下降；若采用常弹性需求 \(Q=A P^\varepsilon\)，在 \(\varepsilon=-0.05\) 且供应缺口较大时会给出极高的机械上界：

\begin{equation}
  P_{\mathrm{ce}}
  =
  P_0
  \left(
    1-\frac{I}{D_0}
  \right)^{1/\varepsilon}.
\end{equation}

本文同时计算 \(P_{\mathrm{ce}}\) 作为稳健性参照，但不将其作为主图结论。两个口径的共同作用，是证明传统低弹性供需框架会显著高估现实价格，而不是精确预测冲突窗口逐日油价。

\subsection{基准模型结果}

\begin{table}[H]
\centering
\caption{传统供需反事实基准模型结果}
\begin{tabular}{lrrrr}
\toprule
情景 & 供应中断量 & 线性需求价格 & 常弹性上界 & 线性倍数 \\
\midrule
低中断 & 1400 & 278.20 & 1494.86 & 2.44 \\
中中断 & 1600 & 307.48 & 2393.23 & 2.70 \\
高中断 & 1800 & 336.77 & 3875.19 & 2.95 \\
\bottomrule
\end{tabular}
\end{table}

\begin{figure}[H]
\centering
\includegraphics[width=0.88\linewidth]{传统供需基准与真实价格对比.png}
\caption{传统供需基准模型与附件真实价格对比。低弹性静态模型明显高估现实价格。}
\end{figure}

结果显示，线性需求反事实基准给出的 278--337 美元/桶显著高于附件冲突窗口最高收盘价 114.06 美元/桶；常弹性机械上界更高，达到 1495--3875 美元/桶。两个口径都说明，单纯依靠供应缺口和低短期弹性无法解释现实价格，必须进一步引入 SPR、库存、绕道运输、需求收缩、价格传导折价和市场预期修复。
